% !TeX root = main.tex
\resetproblem
\begin{center}
    {\LARGE \textbf{Solution Manual of Problem Set 26} \par}
    \vspace{1em}
    {\large Bufan Zheng\\ \href{mailto:bufan@g.ecc.u-tokyo.ac.jp}{\texttt{bufan@g.ecc.u-tokyo.ac.jp}} \par}
\end{center}

\begin{problem}
    \textbf{Scalar in curved space: diff invariance.} For
    \[
        S=\int d^Dx\,\sqrt{-g}\left(\frac{1}{2}g^{\mu\nu}\partial_\mu\varphi\,\partial_\nu\varphi-\frac{m^2}{2}\varphi^2\right),
    \]
    verify that under an infinitesimal diffeomorphism generated by \(\xi^\mu\), the variation is a total derivative. Compute \(\delta\varphi=\xi^\mu\partial_\mu\varphi\) and \(\delta g_{\mu\nu}=\nabla_\mu\xi_\nu+\nabla_\nu\xi_\mu\).
\end{problem}
\begin{solution}[26.1]
	First, let us compute the variation of the Lagrangian:
	\[
	\begin{aligned}\delta\mathcal{L}&=\delta\left(\frac{1}{2}g^{\mu\nu}\nabla_\mu\varphi\nabla_\nu\varphi\right)-\delta\left(\frac{m^2}{2}\varphi^2\right)\\&=\frac{1}{2}(\delta g^{\mu\nu})\nabla_\mu\varphi\nabla_\nu\varphi+g^{\mu\nu}\delta(\nabla_\mu\varphi)\nabla_\nu\varphi-m^2\varphi\delta\varphi\\&=\frac{1}{2}{\left[-(\nabla^\mu\xi^\nu+\nabla^\nu\xi^\mu)\right]}\nabla_\mu\varphi\nabla_\nu\varphi+g^{\mu\nu}\nabla_\mu(\delta\varphi)\nabla_\nu\varphi-m^2\varphi(\xi^\rho\nabla_\rho\varphi)\\&=-(\nabla^\mu\xi^\nu)\nabla_\mu\varphi\nabla_\nu\varphi+g^{\mu\nu}\nabla_\mu(\xi^\rho\nabla_\rho\varphi)\nabla_\nu\varphi+\xi^\rho\nabla_\rho\left(-\frac{m^2}{2}\varphi^2\right)\\&=-(\nabla^\mu\xi^\nu)\nabla_\mu\varphi\nabla_\nu\varphi+g^{\mu\nu}\left[(\nabla_\mu\xi^\rho)\nabla_\rho\varphi+\xi^\rho\nabla_\mu\nabla_\rho\varphi\right]\nabla_\nu\varphi+\xi^\rho\nabla_\rho\left(-\frac{m^2}{2}\varphi^2\right)\\&=-(\nabla^\mu\xi^\nu)\nabla_\mu\varphi\nabla_\nu\varphi+(\nabla^\mu\xi^\rho)\nabla_\rho\varphi\nabla_\mu\varphi+\xi^\rho g^{\mu\nu}\nabla_\mu\nabla_\rho\varphi\nabla_\nu\varphi+\xi^\rho\nabla_\rho\left(-\frac{m^2}{2}\varphi^2\right)\\&=\xi^\rho g^{\mu\nu}\nabla_\mu\nabla_\rho\varphi\nabla_\nu\varphi+\xi^\rho\nabla_\rho\left(-\frac{m^2}{2}\varphi^2\right)\\&=\xi^\rho\nabla_\rho{\left(\frac{1}{2}g^{\mu\nu}\nabla_\mu\varphi\nabla_\nu\varphi\right)}+\xi^\rho\nabla_\rho{\left(-\frac{m^2}{2}\varphi^2\right)}\\&=\xi^\rho\nabla_\rho\mathcal{L}.\end{aligned}
	\]
	Next, the variation of action is:
	\[
	\begin{aligned}\delta S&=\delta\int d^Dx\sqrt{-g}\mathcal{L}=\int d^Dx\delta(\sqrt{-g}\mathcal{L})\\&=\int d^Dx\left[(\delta\sqrt{-g})\mathcal{L}+\sqrt{-g}\left.\delta\mathcal{L}\right]\right.\\&=\int d^Dx{\left[\frac{1}{2}\sqrt{-g}\left.g^{\mu\nu}\delta g_{\mu\nu}\right.\mathcal{L}+\sqrt{-g}\left.\delta\mathcal{L}\right]\right.}\\&=\int d^Dx{\left[\frac{1}{2}\sqrt{-g}\right.}g^{\mu\nu}(\nabla_\mu\xi_\nu+\nabla_\nu\xi_\mu)\left.\mathcal{L}+\sqrt{-g}\right.\delta\mathcal{L}]\\&=\int d^Dx{\left[\sqrt{-g}\left(\nabla_\mu\xi^\mu\right)\mathcal{L}+\sqrt{-g}\delta\mathcal{L}\right]}.\end{aligned}
	\]
	Combining this two results, we obtain:
	\[
	\begin{aligned}\delta S&=\int d^Dx\sqrt{-g}{\left[(\nabla_\mu\xi^\mu)\mathcal{L}+\delta\mathcal{L}\right]}\\&=\int d^Dx\sqrt{-g}{\left[(\nabla_\mu\xi^\mu)\mathcal{L}+\xi^\mu\nabla_\mu\mathcal{L}\right]}\\&=\int d^Dx\sqrt{-g}\nabla_\mu(\xi^\mu\mathcal{L})\\&=\int d^Dx\operatorname{\partial}_\mu\left(\sqrt{-g}\operatorname{\xi}^\mu\mathcal{L}\right).\end{aligned}
	\]
	In the last step, we used $\sqrt{-g}\nabla_\mu V^\mu=\partial_\mu(\sqrt{-g}V^\mu)$. Then we complete the proof.
\end{solution}
\begin{problem}
    \textbf{Background gauge field as spurion.} For a complex scalar with \(D_\mu=\partial_\mu+iA_\mu\), show that if \(\phi\to e^{i\alpha}\phi\) and \(A_\mu\to A_\mu-\partial_\mu\alpha\), then \(|D_\mu\phi|^2\) is invariant.
\end{problem}
\begin{solution}[26.2]
	\[
	\begin{aligned}
			D_\mu^\prime\phi^\prime&=(\partial_\mu+iA_\mu^\prime)(e^{i\alpha}\phi)=\partial_\mu(e^{i\alpha}\phi)+i(A_\mu-\partial_\mu\alpha)e^{i\alpha}\phi\\
		&=e^{i\alpha}(i\partial_\mu\alpha)\phi+e^{i\alpha}\partial_\mu\phi+i(A_\mu-\partial_\mu\alpha)e^{i\alpha}\phi=e^{i\alpha}(\partial_\mu+iA_\mu)\phi=e^{i\alpha}D_\mu\phi.
	\end{aligned}
	\]
	which implies $|D_\mu\phi|^2$ is invariant.
\end{solution}
\begin{problem}
    \textbf{Ward identity with background gauge field.} Consider \(Z[A]=\int D\phi\,D\phi^*\,\exp(iS[\phi;A])\). Derive the Ward identity
    \[
        \nabla_\mu\langle j^\mu(x)\rangle_A=0
    \]
    by varying \(\alpha(x)\) in the background gauge symmetry.
\end{problem}
\begin{solution}[26.3]
	We assume the measurement $D\phi$ and $D\phi^*$ is invariant under $\phi\to e^{i\alpha}\phi$, i.e. this theory is anomaly free. The generate function $Z[A]$ should be invariant under the gauge transformation $\delta A_{\mu}(x)=\nabla_\mu \alpha(x)$:
	\[
	\begin{aligned}
		0=\frac{\delta Z}{\delta A_\mu(x)}&=\int D\phi D\phi^*\frac{\delta}{\delta A_\mu(x)}e^{iS[\phi;A]}\\&=\int D\phi D\phi^*i\underbrace{\left(\frac{\delta S}{\delta A_\mu(x)}\right)}_{j^\mu}e^{iS[\phi;A]}\\
		&=\int d^4x\left(i\langle j^\mu(x)\rangle_AZ[A]\right)\delta A_\mu(x)\\
		&=iZ[A]\int d^4x\langle j^\mu(x)\rangle_A\nabla_\mu\alpha(x)\\
		&\Rightarrow\int\nabla_\mu \braket{j^\mu(x)}_A\alpha(x)=0
		\end{aligned}
	\]
	Since this equation holds for any $\alpha(x)$, we obtain:
	\[
	\boxed{
	\nabla_\mu\langle j^\mu(x)\rangle_A=0
}
	\]
\end{solution}
\begin{problem}
    \textbf{Spurionic mass term.} For a Dirac fermion in curved space,
    \[
        S=\int d^Dx\,\sqrt{-g}\,\bar{\psi}(i\gamma^\mu\nabla_\mu-m)\psi,
    \]
    treat \(m(x)\) as a spurion field. Under a Weyl scaling \(g_{\mu\nu}\to e^{2\sigma}g_{\mu\nu}\), choose a scaling of \(m(x)\) so that the mass term scales homogeneously. What is the Weyl weight of \(m\)?
\end{problem}
\begin{problem}
    \textbf{Holonomy with background \(U(1)\) on \(T^2\).} On \(T^2\) with cycles \(a,b\), let \(A\) be flat but with holonomies \(U_a=\exp(i\oint_a A)\) and \(U_b=\exp(i\oint_b A)\). Show that gauge-invariant data are \((U_a,U_b)\) modulo large gauge transformations that shift the integrals by \(2\pi\mathbb{Z}\).
\end{problem}
\begin{problem}
    \textbf{Flux sector and background metric.} Let \(\Sigma=S^2\) with radius \(R\) and background metric. For a \(U(1)\) background with \(\int_{S^2}F=2\pi k\), take \(F\) proportional to the volume form. Compute \(\int d^2x\,\sqrt{g}\,F^{\mu\nu}F_{\mu\nu}\) in terms of \(k,R\).
\end{problem}
\begin{problem}
    \textbf{Spurionic general covariance and stress tensor.} Couple a QFT to \(g_{\mu\nu}(x)\). Show that diffeomorphism invariance implies the Ward identity
    \[
        \nabla_\mu\langle T^{\mu\nu}\rangle=0
    \]
    in the absence of gravitational anomalies, by varying the generating functional under \(\delta g_{\mu\nu}=\nabla_\mu\xi_\nu+\nabla_\nu\xi_\mu\). If the theory is also coupled to a background \(U(1)\) gauge field \(A_\mu\), show the Ward identity generalizes to
    \[
        \nabla_\mu\langle T^\mu{}_\nu\rangle=\langle j^\mu\rangle F_{\mu\nu}.
    \]
\end{problem}
